% Auto-generated complete chapter from 12_application_system_and_tracker.md
\chapter{12\_application\_system\_and\_tracker.md}
\label{complete:root_014}

\begin{sourcemeta}
\textbf{Fuente:} \href{file:///E:/Laboral/12_application_system_and_tracker.md}{12\_application\_system\_and\_tracker.md}\\
\textbf{Seccion:} root\\
\textbf{Estado:} canonical\\
\textbf{Rol:} Modulo principal\\
\textbf{Lineas originales:} 1327\\
\textbf{Ultima modificacion:} 2026-06-11 14:23\\
\textbf{Deduplicacion conservadora:} 0 bloques repetidos omitidos; 0 caracteres repetidos no reimpresos.
\end{sourcemeta}

\section*{Resumen de orientacion}
This file defines the operating system for Alexander Woodcock Salomón’s international job search.

\section*{Contenido completo depurado}
\addcontentsline{toc}{section}{Contenido completo depurado}




\section{Application System and Tracker}




\subsection{0. Purpose}




This file defines the operating system for Alexander Woodcock Salomón’s international job search.




The goal is not to “apply more.” The goal is to run a controlled pipeline that converts the existing portfolio assets, especially TwinSight X500, into interviews, offers, contractor opportunities, and market feedback.




This system assumes the previous strategic decisions are already fixed:




\begin{itemize}
\item Primary target: Real-Time 3D / Interactive Technical Visualization / Unity WebGL.
\item Secondary target: Unity Technical Artist focused on optimization, runtime systems, tools, shaders, and CAD-to-realtime workflows.
\item Expansion target: Digital Twin / Simulation / XR / Industrial Visualization.
\item Controlled secondary route: Tools / Python Automation / AI-assisted production tooling.
\item Fallback routes: Full-stack and 3D generalist only when strategically useful.
\item Current priority: remote from Colombia.
\item Future mobility: Portuguese passport expected in approximately two years; Germany route possible through legal marriage and residence with German partner.
\end{itemize}




This section is operational. It does not redefine the positioning. It translates the strategy into execution.




\medskip\hrule\medskip




\subsection{1. Core operating principle}




The job search must be treated as a pipeline, not as scattered applications.




Every target should move through a controlled sequence:




\begin{mdcode}
Bloque de codigo: text
Target identified → Fit scored → Evidence selected → Application customized → Outreach sent → Follow-up → Interview → Test → Offer → Decision → Retrospective
\end{mdcode}




The system must prevent four failure modes:




\begin{enumerate}
\item Applying to roles without enough fit.
\item Spending too much time customizing low-probability applications.
\item Over-positioning into roles where the evidence is weak.
\item Losing track of companies, conversations, follow-ups, and feedback.
\end{enumerate}




The tracker is the source of operational truth.




\medskip\hrule\medskip




\subsection{2. Application categories}




All opportunities must be classified before applying.




\subsection{2.1 Category A — High-fit targets}




These are the highest-priority opportunities.




They match TwinSight directly and do not require major repositioning.




Examples:




\begin{itemize}
\item Interactive 3D Developer.
\item Real-Time 3D Developer.
\item Unity WebGL Developer.
\item Technical Visualization Developer.
\item Digital Twin Visualization Developer.
\item Simulation Visualization Developer.
\item Unity Developer for enterprise 3D.
\item XR Developer for industrial or training applications.
\end{itemize}




A Category A opportunity should satisfy most of these conditions:




\begin{itemize}
\item Unity, WebGL, real-time 3D, simulation, digital twin, XR, or technical visualization is central to the role.
\item A technical portfolio case study is relevant.
\item TwinSight can be used as the primary proof.
\item Remote, contractor, LATAM, or global work is plausible.
\item The role does not require senior commercial experience as a hard gate.
\item The company works in industrial, education, simulation, XR, training, product visualization, or enterprise 3D.
\end{itemize}




Application effort: high.




Customization level: strong.




Expected weekly volume: 5–8.




\medskip\hrule\medskip




\subsection{2.2 Category B — Adjacent targets}




These are plausible but require adaptation.




Examples:




\begin{itemize}
\item Unity Technical Artist in games.
\item Game Technical Artist.
\item Shader Technical Artist.
\item Tools Developer.
\item Pipeline Developer.
\item 3D product configurator developer.
\item WebGL / Three.js Developer.
\item Python Automation Developer.
\item AI Tools Developer.
\end{itemize}




A Category B opportunity is valid when:




\begin{itemize}
\item The role is related to the portfolio but not a direct TwinSight match.
\item ARA, Blender portrait, or shader work may support the application.
\item Some missing skills exist, but the gap is not fatal.
\item The application can be made credible without exaggeration.
\end{itemize}




Application effort: medium.




Customization level: moderate.




Expected weekly volume: 8–12.




\medskip\hrule\medskip




\subsection{2.3 Category C — Volume targets}




These are lower-fit but useful for market testing, interview practice, or income fallback.




Examples:




\begin{itemize}
\item Junior Unity Developer.
\item C\# Developer with light Unity.
\item Full-stack developer.
\item Python automation support.
\item Freelance web projects.
\item 3D generalist roles.
\item Simple contractor gigs.
\end{itemize}




A Category C opportunity is acceptable when:




\begin{itemize}
\item It can produce income quickly.
\item It does not consume excessive customization time.
\item It does not force a false professional identity.
\item It can be applied to with a lighter version of the CV.
\end{itemize}




Application effort: low.




Customization level: minimal.




Expected weekly volume: 10–20.




\medskip\hrule\medskip




\subsection{2.4 Watchlist targets}




These are not immediate targets but should be monitored.




Examples:




\begin{itemize}
\item Senior Technical Artist AAA.
\item Senior Simulation Engineer.
\item Digital Twin Architect.
\item Research Engineer.
\item Graphics Programmer.
\item AI Engineer.
\item EU-only roles requiring current EU work authorization.
\item Germany-only roles requiring German language.
\end{itemize}




Watchlist targets should not consume active application time unless a specific opening unusually matches the profile.




\medskip\hrule\medskip




\subsection{3. Weekly application mix}




A balanced weekly target should avoid both over-customization and indiscriminate spam.




Recommended baseline after the thesis defense:




\begin{center}
\scriptsize
\begin{longtable}{>{\raggedright\arraybackslash}p{0.455\linewidth} >{\raggedright\arraybackslash}p{0.455\linewidth}}
\toprule
Category & Weekly target \\
\midrule
\endfirsthead
\toprule
Category & Weekly target \\
\midrule
\endhead
A & 5–8 \\
B & 8–12 \\
C & 10–20 \\
Speculative outreach & 10–15 \\
Recruiter contacts & 5–10 \\
\bottomrule
\end{longtable}
\end{center}




Total weekly outbound actions:




\begin{mdcode}
Bloque de codigo: text
Minimum controlled pace: 30 actions/week
Strong pace: 45–55 actions/week
Aggressive pace: 60–75 actions/week
\end{mdcode}




An outbound action means one of the following:




\begin{itemize}
\item Submitted application.
\item Recruiter message.
\item Founder / technical lead outreach.
\item Referral request.
\item Follow-up.
\item Portfolio feedback request.
\item Community post with job-search relevance.
\end{itemize}




Do not count passive browsing as execution.




\medskip\hrule\medskip




\subsection{4. Daily execution blocks}




The job search should be split into repeatable blocks.




\subsection{4.1 Minimum daily block}




Use this when time or energy is limited.




\begin{mdcode}
Bloque de codigo: text
Duration: 90 minutes

1. Review tracker: 10 min
2. Find opportunities: 20 min
3. Apply or message: 40 min
4. Update tracker: 10 min
5. Next-action planning: 10 min
\end{mdcode}




Expected output:




\begin{itemize}
\item 2–4 applications or messages.
\item Tracker updated.
\item No loose tabs left open.
\end{itemize}




\medskip\hrule\medskip




\subsection{4.2 Standard daily block}




Use this as the default after defense.




\begin{mdcode}
Bloque de codigo: text
Duration: 3 hours

1. Pipeline review: 15 min
2. Target search: 30 min
3. Custom applications: 75 min
4. Outreach/follow-up: 30 min
5. Tracker update: 15 min
6. Portfolio/CV adjustment: 15 min
\end{mdcode}




Expected output:




\begin{itemize}
\item 4–8 outbound actions.
\item 1–2 high-quality customized applications.
\item Follow-ups sent on schedule.
\end{itemize}




\medskip\hrule\medskip




\subsection{4.3 Deep work block}




Use this 2–3 times per week for strategic assets.




\begin{mdcode}
Bloque de codigo: text
Duration: 2–4 hours

Allowed work:
- TwinSight case study refinement.
- Demo video update.
- GitHub README polish.
- Portfolio section improvement.
- Technical test preparation.
- Interview rehearsal.
\end{mdcode}




This block is not for browsing job boards.




\medskip\hrule\medskip




\subsection{5. Weekly schedule}




\subsection{5.1 Recommended weekly structure}




\begin{mdcode}
Bloque de codigo: text
Monday:
- Review metrics.
- Refresh target list.
- Apply to new Category A roles.

Tuesday:
- Category B applications.
- Recruiter outreach.
- Portfolio refinement.

Wednesday:
- Category A applications.
- Speculative outreach to companies.
- Follow-ups.

Thursday:
- Category B/C applications.
- Community networking.
- Interview preparation.

Friday:
- Follow-ups.
- Tracker cleanup.
- Weekly retrospective.

Saturday:
- Optional portfolio or technical skill sprint.

Sunday:
- No active applications unless urgent.
- Light planning only.
\end{mdcode}




\medskip\hrule\medskip




\subsection{5.2 Weekly minimum deliverables}




At the end of each week, the following must exist:




\begin{itemize}
\item All new applications logged.
\item All follow-up dates assigned.
\item All responses categorized.
\item All rejections tagged with reason if known.
\item At least one portfolio or profile asset improved.
\item Weekly metrics written down.
\end{itemize}




If the tracker is not updated, the week is not operationally complete.




\medskip\hrule\medskip




\subsection{6. Tracker architecture}




The tracker should be built in Google Sheets, Notion, Airtable, or a local spreadsheet.




Preferred structure:




\begin{mdcode}
Bloque de codigo: text
1. Dashboard
2. Applications
3. Companies
4. Contacts
5. Follow-ups
6. Interviews
7. Technical tests
8. Offers
9. Metrics
10. Lessons learned
\end{mdcode}




A spreadsheet is sufficient. Notion is optional.




\medskip\hrule\medskip




\subsection{7. Applications table}




\subsection{7.1 Required columns}




Use these fields for every application.




\begin{center}
\scriptsize
\begin{longtable}{>{\raggedright\arraybackslash}p{0.455\linewidth} >{\raggedright\arraybackslash}p{0.455\linewidth}}
\toprule
Field & Values \\
\midrule
\endfirsthead
\toprule
Field & Values \\
\midrule
\endhead
ID & APP-0001 \\
Date applied & YYYY-MM-DD \\
Company & Name \\
Role title & Exact title \\
Category & A / B / C / Watchlist \\
Source & LinkedIn / Job board / Referral \\
Location & Country / Remote \\
Remote scope & WW / LATAM / US / EU / Hybrid \\
Contract type & Employee / Contractor / Unknown \\
Salary range & Stated / Estimated \\
Status & Pipeline status \\
Next action & Specific action \\
Follow-up date & YYYY-MM-DD \\
\bottomrule
\end{longtable}
\end{center}




\medskip\hrule\medskip




\subsection{7.2 Fit columns}




Use these scoring fields to prevent bad applications.




\begin{center}
\scriptsize
\begin{longtable}{>{\raggedright\arraybackslash}p{0.455\linewidth} >{\raggedright\arraybackslash}p{0.455\linewidth}}
\toprule
Field & Scale \\
\midrule
\endfirsthead
\toprule
Field & Scale \\
\midrule
\endhead
Role fit & 1–5 \\
Portfolio fit & 1–5 \\
Remote feasibility & 1–5 \\
Compensation potential & 1–5 \\
Probability & 1–5 \\
Priority & A / B / C \\
\bottomrule
\end{longtable}
\end{center}




Scoring rule:




\begin{mdcode}
Bloque de codigo: text
Total score = Role fit + Portfolio fit + Remote feasibility + Compensation potential + Probability
\end{mdcode}




Interpretation:




\begin{center}
\scriptsize
\begin{longtable}{>{\raggedright\arraybackslash}p{0.455\linewidth} >{\raggedright\arraybackslash}p{0.455\linewidth}}
\toprule
Score & Meaning \\
\midrule
\endfirsthead
\toprule
Score & Meaning \\
\midrule
\endhead
21–25 & Strong target \\
17–20 & Apply selectively \\
13–16 & Low effort only \\
<13 & Do not apply \\
\bottomrule
\end{longtable}
\end{center}




\medskip\hrule\medskip




\subsection{7.3 Evidence columns}




Each application should record which proof was used.




\begin{center}
\scriptsize
\begin{longtable}{>{\raggedright\arraybackslash}p{0.455\linewidth} >{\raggedright\arraybackslash}p{0.455\linewidth}}
\toprule
Field & Examples \\
\midrule
\endfirsthead
\toprule
Field & Examples \\
\midrule
\endhead
Primary proof & TwinSight \\
Secondary proof & ARA / Blender portrait \\
Portfolio link used & URL \\
GitHub link used & URL \\
Demo link used & URL \\
CV variant & Unity / WebGL / Digital Twin \\
Cover note & Yes / No \\
\bottomrule
\end{longtable}
\end{center}




This prevents using the same generic proof for every role.




\medskip\hrule\medskip




\subsection{8. Company table}




The company table prevents duplicate research.




Required fields:




\begin{center}
\scriptsize
\begin{longtable}{>{\raggedright\arraybackslash}p{0.455\linewidth} >{\raggedright\arraybackslash}p{0.455\linewidth}}
\toprule
Field & Values \\
\midrule
\endfirsthead
\toprule
Field & Values \\
\midrule
\endhead
Company & Name \\
Website & URL \\
Careers page & URL \\
LinkedIn & URL \\
Category & Digital Twin / XR / WebGL \\
Region & US / EU / LATAM \\
Language & English / German / Portuguese \\
Remote scope & WW / LATAM / EU / Unknown \\
Contractor likelihood & High / Medium / Low \\
Priority & A / B / C / Watchlist \\
Last checked & YYYY-MM-DD \\
\bottomrule
\end{longtable}
\end{center}




Additional notes:




\begin{itemize}
\item Keep company notes short.
\item Do not research the same company repeatedly.
\item If no suitable opening exists, mark it as \texttt{Watchlist}.
\item Review Watchlist companies every 30 days.
\end{itemize}




\medskip\hrule\medskip




\subsection{9. Contacts table}




Any person contacted must be logged.




Required fields:




\begin{center}
\scriptsize
\begin{longtable}{>{\raggedright\arraybackslash}p{0.455\linewidth} >{\raggedright\arraybackslash}p{0.455\linewidth}}
\toprule
Field & Values \\
\midrule
\endfirsthead
\toprule
Field & Values \\
\midrule
\endhead
Name & Contact name \\
Role & Recruiter / Lead / Founder \\
Company & Company \\
LinkedIn & URL \\
Email & If public \\
Relationship & Cold / Warm / Referral \\
Last contact & YYYY-MM-DD \\
Next follow-up & YYYY-MM-DD \\
Status & No reply / Replied / Call \\
\bottomrule
\end{longtable}
\end{center}




Rules:




\begin{itemize}
\item Do not message the same person repeatedly without a reason.
\item Follow up once after 5–7 business days.
\item A second follow-up is only justified for high-fit targets.
\item No third follow-up unless they previously replied.
\end{itemize}




\medskip\hrule\medskip




\subsection{10. Status pipeline}




Use a controlled status list.




\begin{mdcode}
Bloque de codigo: text
Identified
Scored
Applied
Outreach sent
Follow-up due
Follow-up sent
Recruiter screen
Technical interview
Portfolio review
Technical test
Final interview
Offer
Rejected
Ghosted
Withdrawn
Archived
\end{mdcode}




Do not create custom statuses for every case. Too many statuses make the tracker unusable.




\medskip\hrule\medskip




\subsection{11. Follow-up rules}




\subsection{11.1 Application follow-up}




For ordinary applications:




\begin{mdcode}
Bloque de codigo: text
Day 0: Apply
Day 7: Follow up if contact exists
Day 14: Final follow-up if high priority
Day 21: Archive if no response
\end{mdcode}




For high-fit Category A applications:




\begin{mdcode}
Bloque de codigo: text
Day 0: Apply + targeted outreach
Day 5–7: Follow up
Day 14: Second follow-up or new contact
Day 30: Archive or move to Watchlist
\end{mdcode}




\medskip\hrule\medskip




\subsection{11.2 Recruiter follow-up}




\begin{mdcode}
Bloque de codigo: text
Day 0: Contact recruiter
Day 5–7: Follow up
Day 14: Archive if no answer
\end{mdcode}




If a recruiter replies but has no current role, set reminder for 30–45 days.




\medskip\hrule\medskip




\subsection{11.3 Interview follow-up}




\begin{mdcode}
Bloque de codigo: text
Same day: thank-you note
Day 5: follow up if no timeline was given
Day 10: follow up if timeline expired
\end{mdcode}




Do not over-message after interviews.




\medskip\hrule\medskip




\subsection{12. Customization rules}




Not every application deserves the same effort.




\subsection{12.1 Category A customization}




Category A applications require:




\begin{itemize}
\item CV variant selected intentionally.
\item Short customized opening message.
\item TwinSight case study link.
\item Relevant bullet reordered near the top.
\item Company-specific reason.
\item Follow-up contact identified.
\end{itemize}




Expected time per application:




\begin{mdcode}
Bloque de codigo: text
30–60 minutes
\end{mdcode}




\medskip\hrule\medskip




\subsection{12.2 Category B customization}




Category B applications require:




\begin{itemize}
\item Correct CV variant.
\item One paragraph tailored to the role.
\item Relevant project link.
\item Basic tracker entry.
\end{itemize}




Expected time per application:




\begin{mdcode}
Bloque de codigo: text
15–30 minutes
\end{mdcode}




\medskip\hrule\medskip




\subsection{12.3 Category C customization}




Category C applications require:




\begin{itemize}
\item Correct CV variant.
\item No long cover letter.
\item No excessive research.
\item Quick tracker entry.
\end{itemize}




Expected time per application:




\begin{mdcode}
Bloque de codigo: text
5–15 minutes
\end{mdcode}




Category C should not steal time from portfolio and Category A roles.




\medskip\hrule\medskip




\subsection{13. CV variant selection}




Use the correct CV variant for each role.




\begin{center}
\scriptsize
\begin{longtable}{>{\raggedright\arraybackslash}p{0.455\linewidth} >{\raggedright\arraybackslash}p{0.455\linewidth}}
\toprule
Role type & CV variant \\
\midrule
\endfirsthead
\toprule
Role type & CV variant \\
\midrule
\endhead
Interactive 3D & Real-Time 3D \\
Unity WebGL & WebGL / Interactive 3D \\
Digital Twin & Technical Visualization \\
Simulation & Technical Visualization \\
XR & XR / Unity \\
Technical Artist & Unity Technical Artist \\
Tools & Tools / Python \\
AI automation & Python / AI Tools \\
Full-stack & Full-stack fallback \\
3D generalist & 3D / Technical Art fallback \\
\bottomrule
\end{longtable}
\end{center}




Do not send the AI Tools CV to Unity roles unless the job explicitly values automation.




Do not send a 3D Artist CV to technical visualization roles.




\medskip\hrule\medskip




\subsection{14. Portfolio link selection}




Each application should include only the most relevant links.




\subsection{14.1 Default link order}




For most roles:




\begin{mdcode}
Bloque de codigo: text
1. Portfolio homepage
2. TwinSight case study
3. GitHub profile
4. LinkedIn
\end{mdcode}




For technical roles:




\begin{mdcode}
Bloque de codigo: text
1. TwinSight case study
2. GitHub repository
3. Demo video
4. Portfolio homepage
\end{mdcode}




For tools / automation roles:




\begin{mdcode}
Bloque de codigo: text
1. ARA case study
2. GitHub repository
3. TwinSight as secondary proof
4. LinkedIn
\end{mdcode}




For art-adjacent roles:




\begin{mdcode}
Bloque de codigo: text
1. TwinSight case study
2. Blender portrait technical breakdown
3. ArtStation
4. Portfolio homepage
\end{mdcode}




\medskip\hrule\medskip




\subsection{15. Application quality gates}




Before applying, check these gates.




\subsection{15.1 Mandatory gate}




Do not apply if all are true:




\begin{itemize}
\item Requires current US/EU work authorization.
\item Does not allow remote or contractor work.
\item Requires 5+ years formal industry experience.
\item Has no connection to Unity, 3D, simulation, XR, WebGL, tools, or Python.
\item Would require a false claim.
\end{itemize}




\medskip\hrule\medskip




\subsection{15.2 Apply anyway gate}




Apply despite imperfect fit if at least three are true:




\begin{itemize}
\item TwinSight is directly relevant.
\item Remote scope is worldwide, LATAM, or contractor-friendly.
\item Requirements are flexible.
\item Role is junior, mid-junior, associate, contractor, or project-based.
\item Company works in simulation, digital twin, XR, education, product visualization, or industrial 3D.
\item The application can be customized in less than 45 minutes.
\end{itemize}




\medskip\hrule\medskip




\subsection{16. Outreach system}




This section defines cadence only. Full outreach scripts belong in a later file.




\subsection{16.1 Who to contact}




Priority contacts:




\begin{enumerate}
\item Technical Artist Lead.
\item Unity Lead.
\item 3D / XR / Simulation Lead.
\item Founder of small studio.
\item Technical recruiter.
\item Creative technology director.
\item Engineering manager.
\item Former UNAD / UNAL / project contacts.
\end{enumerate}




Avoid contacting random HR profiles without relevance when better technical contacts exist.




\medskip\hrule\medskip




\subsection{16.2 Outreach pairing}




Use this pattern:




\begin{mdcode}
Bloque de codigo: text
Application + targeted message
\end{mdcode}




Do not rely only on job portals.




For Category A roles, always attempt at least one direct contact.




For Category B roles, contact if the company is small or the role is highly relevant.




For Category C roles, direct outreach is optional.




\medskip\hrule\medskip




\subsection{17. Metrics dashboard}




Track weekly metrics.




\subsection{17.1 Activity metrics}




\begin{center}
\scriptsize
\begin{longtable}{>{\raggedright\arraybackslash}p{0.455\linewidth} >{\raggedright\arraybackslash}p{0.455\linewidth}}
\toprule
Metric & Target \\
\midrule
\endfirsthead
\toprule
Metric & Target \\
\midrule
\endhead
Applications sent & 20–40 \\
Category A applications & 5–8 \\
Direct outreach messages & 10–15 \\
Follow-ups sent & 5–15 \\
Recruiters contacted & 5–10 \\
Portfolio improvements & 1–2 \\
\bottomrule
\end{longtable}
\end{center}




\medskip\hrule\medskip




\subsection{17.2 Conversion metrics}




\begin{center}
\scriptsize
\begin{longtable}{>{\raggedright\arraybackslash}p{0.455\linewidth} >{\raggedright\arraybackslash}p{0.455\linewidth}}
\toprule
Metric & Good signal \\
\midrule
\endfirsthead
\toprule
Metric & Good signal \\
\midrule
\endhead
Response rate & >8\% \\
Recruiter screen rate & >3\% \\
Interview rate & >2\% \\
Technical test rate & >1\% \\
Offer rate & >0.3\% \\
\bottomrule
\end{longtable}
\end{center}




These numbers are directional. They are not guarantees.




If response rate is below 3\% after 100 applications, the issue is likely one of these:




\begin{itemize}
\item Positioning is too broad.
\item CV is weak.
\item Portfolio is not accessible enough.
\item Wrong roles are being targeted.
\item Applications are going to local-only jobs.
\item Outreach is missing.
\end{itemize}




\medskip\hrule\medskip




\subsection{18. Weekly review}




Every Friday, complete a short review.




\subsection{18.1 Review questions}




\begin{mdcode}
Bloque de codigo: text
1. How many applications were sent?
2. How many were Category A?
3. How many direct messages were sent?
4. How many responses arrived?
5. Which roles responded?
6. Which roles ignored the application?
7. Which portfolio link was used most?
8. Which objections appeared?
9. What should change next week?
10. What should be stopped?
\end{mdcode}




\medskip\hrule\medskip




\subsection{18.2 Weekly decision rules}




If there are no responses after 50 applications:




\begin{itemize}
\item Review CV headline.
\item Review portfolio homepage.
\item Add demo video thumbnail.
\item Improve TwinSight summary.
\item Increase direct outreach.
\end{itemize}




If there are responses but no interviews:




\begin{itemize}
\item Check salary expectation.
\item Check work authorization wording.
\item Check whether the role requires experience not shown.
\item Improve project explanation.
\end{itemize}




If there are interviews but no progress:




\begin{itemize}
\item Improve verbal explanation.
\item Prepare technical breakdown.
\item Practice AI-assisted development disclosure.
\item Prepare code and architecture discussion.
\end{itemize}




If there are technical tests but poor performance:




\begin{itemize}
\item Create a skill sprint.
\item Reduce application volume temporarily.
\item Focus on Unity/C\#/WebGL fundamentals.
\end{itemize}




\medskip\hrule\medskip




\subsection{19. Monthly review}




At the end of every month, classify the job search.




\subsection{19.1 Outcome categories}




\begin{mdcode}
Bloque de codigo: text
Healthy pipeline:
- Responses arriving.
- Some recruiter screens.
- Some interviews.
- No major positioning contradiction.

Weak pipeline:
- Low responses.
- Few interviews.
- Applications mostly ignored.

Wrong market:
- Roles require US/EU authorization.
- Remote is not actually worldwide.
- Seniority mismatch is frequent.

Wrong evidence:
- Recruiters ask for commercial experience.
- Portfolio does not show enough proof.
- TwinSight not immediately understood.
\end{mdcode}




\medskip\hrule\medskip




\subsection{19.2 Monthly adjustment rules}




After 30 days:




\begin{itemize}
\item Keep applying if response rate is acceptable.
\item If no response, revise CV and portfolio front page.
\item If roles reject due to authorization, shift toward LATAM contractor and global contractor platforms.
\item If roles reject due to experience, increase project-based outreach.
\item If roles reject due to skill gaps, schedule a focused sprint.
\end{itemize}




After 60 days:




\begin{itemize}
\item Double down on the top two responding role families.
\item Stop applying to low-response categories.
\item Use feedback to rewrite positioning.
\item Increase direct outreach and referral attempts.
\end{itemize}




After 90 days:




\begin{itemize}
\item Reassess salary target.
\item Consider contractor/freelance bridge work.
\item Consider temporary fallback route if no traction.
\item Decide whether to invest more in ARA, WebGL, XR, or Unity Technical Art.
\end{itemize}




\medskip\hrule\medskip




\subsection{20. Bad-offer rejection criteria}




Reject or negotiate offers with these red flags.




\subsection{20.1 Structural red flags}




\begin{itemize}
\item No written contract.
\item Unclear payment terms.
\item Payment only after vague milestones.
\item Full-time exclusivity for low contractor pay.
\item Fixed schedule but no employee benefits.
\item Required unpaid test longer than 4–6 hours.
\item No clarity on IP ownership.
\item No clarity on termination terms.
\item Salary below minimum target without learning value.
\end{itemize}




\medskip\hrule\medskip




\subsection{20.2 Remote-work red flags}




\begin{itemize}
\item “Remote” but requires local work authorization.
\item “Remote” but demands relocation soon without support.
\item “Contractor” but requires employee-like control.
\item No mention of payment currency.
\item No invoicing process.
\item No tax documentation.
\item No equipment clarity.
\end{itemize}




\medskip\hrule\medskip




\subsection{20.3 Portfolio/IP red flags}




\begin{itemize}
\item Requests full project files before contract.
\item Asks to recreate TwinSight-like work for free.
\item Requires confidential unpaid work.
\item Claims ownership of test work without compensation.
\item Requests manufacturer/CAD assets without legal clarity.
\end{itemize}




\medskip\hrule\medskip




\subsection{21. Offer evaluation framework}




Evaluate every offer using six dimensions.




\begin{center}
\scriptsize
\begin{longtable}{>{\raggedright\arraybackslash}p{0.455\linewidth} >{\raggedright\arraybackslash}p{0.455\linewidth}}
\toprule
Dimension & Weight \\
\midrule
\endfirsthead
\toprule
Dimension & Weight \\
\midrule
\endhead
Net income & 30\% \\
Role fit & 20\% \\
Portfolio leverage & 15\% \\
International leverage & 15\% \\
Stability & 10\% \\
Learning value & 10\% \\
\bottomrule
\end{longtable}
\end{center}




Use this scoring:




\begin{mdcode}
Bloque de codigo: text
5 = excellent
4 = strong
3 = acceptable
2 = weak
1 = unacceptable
\end{mdcode}




Any offer below 3/5 in both net income and role fit should be rejected unless there is an urgent financial need.




\medskip\hrule\medskip




\subsection{22. Gross-to-net tracking}




Track compensation as gross and estimated net.




Required fields:




\begin{center}
\scriptsize
\begin{longtable}{>{\raggedright\arraybackslash}p{0.455\linewidth} >{\raggedright\arraybackslash}p{0.455\linewidth}}
\toprule
Field & Example \\
\midrule
\endfirsthead
\toprule
Field & Example \\
\midrule
\endhead
Gross monthly & USD 3,000 \\
Payment currency & USD \\
Platform fee & 0–5\% \\
Conversion spread & 0–3\% \\
Social security reserve & Estimate \\
Tax reserve & Estimate \\
Net estimate & Conservative \\
\bottomrule
\end{longtable}
\end{center}




Do not compare offers only by gross salary.




A USD 3,000 contractor offer may be weaker than a lower employee offer if it has unstable payments, no benefits, and high compliance cost.




\medskip\hrule\medskip




\subsection{23. Tracker templates}




\subsection{23.1 Applications CSV schema}




\begin{mdcode}
Bloque de codigo: csv
id,date_applied,company,role_title,category,source,location,remote_scope,contract_type,salary_range,status,next_action,follow_up_date,role_fit,portfolio_fit,remote_feasibility,compensation_potential,probability,total_score,primary_proof,secondary_proof,cv_variant,portfolio_link,github_link,contact_name,contact_url,notes
\end{mdcode}




\medskip\hrule\medskip




\subsection{23.2 Companies CSV schema}




\begin{mdcode}
Bloque de codigo: csv
company,website,careers_url,linkedin_url,category,region,language,remote_scope,contractor_likelihood,priority,last_checked,notes
\end{mdcode}




\medskip\hrule\medskip




\subsection{23.3 Contacts CSV schema}




\begin{mdcode}
Bloque de codigo: csv
name,role,company,linkedin,email,relationship,last_contact,next_follow_up,status,notes
\end{mdcode}




\medskip\hrule\medskip




\subsection{23.4 Interviews CSV schema}




\begin{mdcode}
Bloque de codigo: csv
company,role,interview_stage,date,interviewer,next_step,deadline,questions_asked,weak_points,follow_up_sent,notes
\end{mdcode}




\medskip\hrule\medskip




\subsection{24. File and folder system}




Create a job-search folder with this structure:




\begin{mdcode}
Bloque de codigo: text
/job_search/
  /cv/
    alexander_woodcock_realtime_3d_cv.pdf
    alexander_woodcock_unity_technical_artist_cv.pdf
    alexander_woodcock_webgl_interactive_3d_cv.pdf
    alexander_woodcock_ai_tools_cv.pdf
  /cover_notes/
  /applications/
  /interviews/
  /technical_tests/
  /company_research/
  /portfolio_assets/
  /screenshots/
  /offers/
\end{mdcode}




Never overwrite submitted CVs. Keep dated versions.




Example:




\begin{mdcode}
Bloque de codigo: text
alexander_woodcock_realtime_3d_cv_2026-06-15.pdf
\end{mdcode}




\medskip\hrule\medskip




\subsection{25. Naming conventions}




Use consistent names.




\begin{mdcode}
Bloque de codigo: text
Company_Role_Date_Status
\end{mdcode}




Example:




\begin{mdcode}
Bloque de codigo: text
Treeview_UnityDeveloper_2026-06-18_Applied
\end{mdcode}




For interview notes:




\begin{mdcode}
Bloque de codigo: text
Company_Role_InterviewStage_Date.md
\end{mdcode}




Example:




\begin{mdcode}
Bloque de codigo: text
Treeview_UnityDeveloper_RecruiterScreen_2026-06-24.md
\end{mdcode}




\medskip\hrule\medskip




\subsection{26. Interview feedback capture}




After every interview, write notes within 30 minutes.




Template:




\begin{mdcode}
Bloque de codigo: text
Company:
Role:
Date:
Stage:
Interviewers:

What they asked:

What I answered well:

What I answered poorly:

Technical gaps exposed:

Portfolio questions:

AI-assisted development questions:

Work authorization questions:

Salary questions:

Next action:

Adjustment needed:
\end{mdcode}




This creates feedback loops.




\medskip\hrule\medskip




\subsection{27. Technical test handling}




Before accepting a test, record:




\begin{itemize}
\item Expected time.
\item Deadline.
\item Compensation.
\item IP ownership.
\item Scope.
\item Tools required.
\item Evaluation criteria.
\end{itemize}




Reject or negotiate unpaid tests that require excessive time or production-quality deliverables.




Suggested rule:




\begin{mdcode}
Bloque de codigo: text
Accept unpaid test: ≤4 hours
Question unpaid test: 4–8 hours
Reject or negotiate: >8 hours
\end{mdcode}




Exception: a high-fit, high-salary Category A opportunity may justify a longer test if the scope is clearly bounded.




\medskip\hrule\medskip




\subsection{28. Market feedback tags}




Tag every rejection or weak response.




Useful tags:




\begin{mdcode}
Bloque de codigo: text
No EU authorization
No US authorization
Too junior
Too senior
Needs shipped game
Needs commercial Unity
Needs WebGL proof
Needs shader proof
Needs C++
Needs Unreal
Needs German
Needs portfolio polish
Salary mismatch
No response
Role closed
\end{mdcode}




After 30–50 responses or rejections, these tags reveal the actual bottleneck.




\medskip\hrule\medskip




\subsection{29. Skill sprint trigger system}




Do not start new courses randomly.




Start a skill sprint only when the tracker shows repeated evidence.




\begin{center}
\scriptsize
\begin{longtable}{>{\raggedright\arraybackslash}p{0.455\linewidth} >{\raggedright\arraybackslash}p{0.455\linewidth}}
\toprule
Trigger & Sprint \\
\midrule
\endfirsthead
\toprule
Trigger & Sprint \\
\midrule
\endhead
Needs WebGL proof & Unity WebGL build polish \\
Needs shader proof & Shader Graph/HLSL mini-breakdown \\
Needs tools proof & Unity editor tool demo \\
Needs XR proof & Small XR prototype \\
Needs Python proof & ARA cleanup \\
Needs German & A1/A2 German routine \\
Needs commercial polish & TwinSight case study upgrade \\
\bottomrule
\end{longtable}
\end{center}




A skill sprint should last 1–2 weeks and produce a visible artifact.




\medskip\hrule\medskip




\subsection{30. Anti-dispersion rules}




The main operational risk is not lack of opportunities. It is fragmentation.




Rules:




\begin{enumerate}
\item Do not create a new portfolio project unless it supports an observed market gap.
\item Do not rewrite LinkedIn every week.
\item Do not change positioning after a few rejections.
\item Do not apply to roles that require false claims.
\item Do not study German, Houdini, AI, WebGL, React, and Unreal at the same time.
\item Do not polish low-priority projects before TwinSight is publicly strong.
\item Do not treat browsing as work.
\item Do not measure progress by tabs opened.
\end{enumerate}




\medskip\hrule\medskip




\subsection{31. Execution phases}




\subsection{31.1 Phase 1 — Setup}




Duration:




\begin{mdcode}
Bloque de codigo: text
1–2 weeks
\end{mdcode}




Deliverables:




\begin{itemize}
\item Final CV variants.
\item LinkedIn updated.
\item GitHub pinned repos fixed.
\item TwinSight case study published.
\item Demo video uploaded.
\item Tracker created.
\item Company list imported.
\item First 30 targets scored.
\end{itemize}




\medskip\hrule\medskip




\subsection{31.2 Phase 2 — Controlled launch}




Duration:




\begin{mdcode}
Bloque de codigo: text
Weeks 3–6
\end{mdcode}




Actions:




\begin{itemize}
\item Apply to 20–40 roles per week.
\item Prioritize Category A and B.
\item Contact recruiters.
\item Run follow-up system.
\item Record responses.
\item Update assets based on objections.
\end{itemize}




\medskip\hrule\medskip




\subsection{31.3 Phase 3 — Optimization}




Duration:




\begin{mdcode}
Bloque de codigo: text
Weeks 7–12
\end{mdcode}




Actions:




\begin{itemize}
\item Stop low-response categories.
\item Double down on responding markets.
\item Improve weak proof points.
\item Increase direct outreach.
\item Prepare negotiation.
\item Consider fallback contractor work if needed.
\end{itemize}




\medskip\hrule\medskip




\subsection{32. Minimum viable tracker setup}




If time is limited, create only one sheet with these columns:




\begin{mdcode}
Bloque de codigo: csv
Date,Company,Role,Category,Source,Remote Scope,Contract Type,Salary,Status,Follow-up Date,Fit Score,Portfolio Link,CV Variant,Contact,Notes
\end{mdcode}




This is enough to start.




Do not delay applications because the tracker is too complex.




\medskip\hrule\medskip




\subsection{33. Quality control checklist}




Before each weekly review, check:




\begin{itemize}
\item No duplicate applications.
\item No missing follow-up dates.
\item No applications without category.
\item No high-fit company without direct outreach attempt.
\item No untracked recruiter conversation.
\item No interview without notes.
\item No offer without gross-to-net estimate.
\item No technical test without scope evaluation.
\end{itemize}




\medskip\hrule\medskip




\subsection{34. Final operational recommendation}




The initial target is not to find the perfect job. The initial target is to generate enough market signal to know which version of the profile is converting.




The first 100 outbound actions should be treated as a calibration period.




After 100 outbound actions, evaluate:




\begin{mdcode}
Bloque de codigo: text
1. Which role family responds?
2. Which market responds?
3. Which salary range responds?
4. Which project gets attention?
5. Which requirement blocks progress?
6. Which claim needs adjustment?
7. Which skill gap appears repeatedly?
\end{mdcode}




Then adjust the strategy based on evidence, not intuition.




\medskip\hrule\medskip




\subsection{35. Next file dependency}




This file feeds directly into:




\begin{mdcode}
Bloque de codigo: text
13_interview_negotiation_and_offer_risk.md
14_30_60_90_execution_plan.md
15_final_integrated_strategy.md
\end{mdcode}




The tracker defines the execution loop. The interview and negotiation file will define how to convert interviews into offers and how to reject bad opportunities.



